\documentclass[12pt,a4paper]{article}

\usepackage[a4paper,margin=1in]{geometry}
\usepackage[T1]{fontenc}
\usepackage{lmodern}
\usepackage{xcolor}
\usepackage{eso-pic}
\usepackage{graphicx}
\usepackage{listings}
\usepackage{fancyhdr}
\usepackage{enumitem}

\newcommand{\addwatermark}{%
    \AddToShipoutPictureBG{%
        \AtPageLowerRight{%
            \raisebox{1.2cm}{%
                \makebox[0pt][r]{%
                    \rotatebox{45}{%
                        \textcolor{gray!20}{\fontsize{12}{14}\selectfont
                        Varun B Karmikanda\ --\ 261100610029}%
                    }%
                }%
            }%
        }%
        \AtPageUpperLeft{%
            \raisebox{-1.2cm}{%
                \hspace{1.5cm}\textcolor{white}{Varun B Karmikanda\ --\ 261100610029}%
            }%
        }%
    }%
}
\addwatermark

\definecolor{codegray}{gray}{0.95}
\lstset{
    language=C,
    basicstyle=\ttfamily\scriptsize,
    backgroundcolor=\color{codegray},
    frame=single,
    breaklines=true,
    showstringspaces=false,
    numbers=left,
    numberstyle=\tiny\color{gray},
    numbersep=8pt
}

\pagestyle{fancy}
\fancyhf{}
\lhead{RTOS LAB}
\rhead{Assignment 3}
\cfoot{\thepage}

\newcommand{\outputimage}[1]{%
    \begin{center}
    \includegraphics[width=0.9\textwidth,keepaspectratio]{images/#1}
    \end{center}
}

\begin{document}

\begin{titlepage}
    \centering
    \vspace*{2cm}
    {\Large\textbf{REAL TIME OPERATING SYSTEM}}\\[0.7cm]
    {\Large\textbf{RTOS LAB}}\\[1.2cm]
    {\Huge\textbf{Assignment 3}}\\[2cm]
    \begin{tabular}{r p{7.5cm}}
    \Large{Name:}     & {\Large\textbf{Varun B Karmikanda}} \\[0.4cm]
    \Large{USN:}      & {\Large{261100610029}} \\[0.4cm]
    \textbf{Program:} & ME Computer Science \& \\[0.1cm]
                      & Engineering
\end{tabular}
    \vfill
    {\large \today}
\end{titlepage}

\section*{1. Shared pthread}
\textbf{Question:} Create two threads that access shared and local variables,
and display their process and thread information.

\textbf{Code:}
\begin{lstlisting}
#include <pthread.h>
#include <stdio.h>
#include <unistd.h>

int shared = 100;

static void *worker(void *arg)
{
    int mine = 1;
    pid_t pid = getpid();
    pthread_t me = pthread_self();

    (void)arg;

    printf("pid=%d thread=%lu\n", (int)pid, (unsigned long)me);
    printf("&shared=%p &mine=%p\n",
           (void *)&shared, (void *)&mine);

    sleep(1);
    return NULL;
}

int main(void)
{
    pthread_t t1, t2;

    pthread_create(&t1, NULL, worker, NULL);
    pthread_create(&t2, NULL, worker, NULL);

    pthread_join(t1, NULL);
    pthread_join(t2, NULL);

    return 0;
}
\end{lstlisting}

\outputimage{01-shared-pthread.png}

\newpage
\section*{2. Basic Fork}
\textbf{Question:} Create a child process using \texttt{fork()} and display
the parent and child process IDs.

\textbf{Code:}
\begin{lstlisting}
#include <stdio.h>
#include <unistd.h>
#include <sys/wait.h>

int main() {
    pid_t pid = fork();

    if (pid < 0) {
        perror("fork failed");
        return 1;
    } else if (pid == 0) {
        // Child process
        printf("Child: my pid = %d, parent pid = %d\n", getpid(), getppid());
    } else {
        // Parent process
        printf("Parent: my pid = %d, child pid = %d\n", getpid(), pid);
        wait(NULL);  // wait for child to finish
        printf("Parent: child has exited\n");
    }

    return 0;
}
\end{lstlisting}

\outputimage{02-basic-fork.png}

\newpage
\section*{3. Fork With Sleep}
\textbf{Question:} Create a child process that sleeps for 20 seconds while
the parent waits for it using \texttt{waitpid()}.

\textbf{Code:}
\begin{lstlisting}
#include <stdio.h>
#include <unistd.h>
#include <sys/wait.h>

int main() {
    pid_t pid = fork();

    if (pid < 0) {
        perror("fork failed");
        return 1;
    } else if (pid == 0) {
        // Child process
        printf("Child: my pid = %d, parent pid = %d\n", getpid(), getppid());
        printf("Child: sleeping for 20 seconds...\n");
        fflush(stdout);
        sleep(20);
        printf("Child: woke up, exiting now\n");
    } else {
        // Parent process
        printf("Parent: my pid = %d, child pid = %d\n", getpid(), pid);
        printf("Parent: waiting for child to finish...\n");
        fflush(stdout);

        int status;
        pid_t exited_pid = waitpid(pid, &status, 0);  // blocks here for ~20 sec

        if (exited_pid == pid) {
            if (WIFEXITED(status))
                printf("Parent: child %d exited with status %d\n", exited_pid, WEXITSTATUS(status));
            else
                printf("Parent: child %d did not exit normally\n", exited_pid);
        }
    }

    return 0;
}
\end{lstlisting}

\outputimage{03-fork-sleep.png}

\newpage
\section*{4. Fork Three Children}
\textbf{Question:} Create three child processes with different sleep
durations and wait for all of them.

\textbf{Code:}
\begin{lstlisting}
#include <stdio.h>
#include <unistd.h>
#include <sys/wait.h>

int main() {
    int sleep_times[3] = {10, 20, 30};
    pid_t child_pids[3];

    for (int i = 0; i < 3; i++) {
        pid_t pid = fork();

        if (pid < 0) {
            perror("fork failed");
            return 1;
        } else if (pid == 0) {
            // Child process
            printf("Child %d: pid = %d, sleeping for %d seconds\n",
                   i, getpid(), sleep_times[i]);
            fflush(stdout);
            sleep(sleep_times[i]);
            printf("Child %d: pid = %d, woke up after %d seconds\n",
                   i, getpid(), sleep_times[i]);
            return 0;  // child exits here, doesn't continue the loop
        } else {
            // Parent process: record child's pid and move on to fork the next one
            child_pids[i] = pid;
        }
    }

    // Parent waits for all 3 children
    printf("Parent: pid = %d, waiting for all children...\n", getpid());
    fflush(stdout);

    for (int i = 0; i < 3; i++) {
        int status;
        pid_t exited_pid = waitpid(child_pids[i], &status, 0);
        if (exited_pid == child_pids[i] && WIFEXITED(status)) {
            printf("Parent: child %d (pid %d) exited with status %d\n",
                   i, exited_pid, WEXITSTATUS(status));
        }
    }

    printf("Parent: all children have exited\n");
    return 0;
}
\end{lstlisting}

\outputimage{04-fork-three-children.png}

\newpage
\section*{5. Fork and waitpid}
\textbf{Question:} Use \texttt{waitpid()} to wait specifically for the
child created by the parent.

\textbf{Code:}
\begin{lstlisting}
#include <stdio.h>
#include <unistd.h>
#include <sys/wait.h>

int main() {
    pid_t pid = fork();

    if (pid < 0) {
        perror("fork failed");
        return 1;
    } else if (pid == 0) {
        // Child process
        printf("Child: my pid = %d, parent pid = %d\n", getpid(), getppid());
    } else {
        // Parent process
        printf("Parent: my pid = %d, child pid = %d\n", getpid(), pid);

        int status;
        pid_t exited_pid = waitpid(pid, &status, 0);  // wait specifically for this child's pid

        if (exited_pid == pid) {
            if (WIFEXITED(status))
                printf("Parent: child %d exited with status %d\n", exited_pid, WEXITSTATUS(status));
            else
                printf("Parent: child %d did not exit normally\n", exited_pid);
        }
    }

    return 0;
}
\end{lstlisting}

\outputimage{05-fork-waitpid.png}

\newpage
\section*{6. Recursive Fork}
\textbf{Question:} Create a chain of processes where each child creates
another child, and use \texttt{waitpid()} between parent and child processes.

\textbf{Code:}
\begin{lstlisting}
#include <stdio.h>
#include <unistd.h>
#include <sys/wait.h>

int main() {
    printf("Original parent: pid = %d\n", getpid());
    fflush(stdout);

    pid_t pid1 = fork();

    if (pid1 < 0) {
        perror("fork failed");
        return 1;
    } else if (pid1 == 0) {
        // ---- child1 ----
        printf("Child1: pid = %d, parent pid = %d\n", getpid(), getppid());
        fflush(stdout);

        pid_t pid2 = fork();

        if (pid2 < 0) {
            perror("fork failed");
            return 1;
        } else if (pid2 == 0) {
            // ---- child2 ----
            printf("Child2: pid = %d, parent pid = %d\n", getpid(), getppid());
            fflush(stdout);

            pid_t pid3 = fork();

            if (pid3 < 0) {
                perror("fork failed");
                return 1;
            } else if (pid3 == 0) {
                // ---- child3 ----
                printf("Child3: pid = %d, parent pid = %d, sleeping 10s\n", getpid(), getppid());
                fflush(stdout);
                sleep(10);
                printf("Child3: pid = %d, woke up, exiting\n", getpid());
                return 0;
            } else {
                // child2 waits for child3
                int status;
                waitpid(pid3, &status, 0);
                printf("Child2: pid = %d, child3 (pid %d) exited\n", getpid(), pid3);
                sleep(10);
                printf("Child2: pid = %d, woke up, exiting\n", getpid());
                return 0;
            }
        } else {
            // child1 waits for child2
            int status;
            waitpid(pid2, &status, 0);
            printf("Child1: pid = %d, child2 (pid %d) exited\n", getpid(), pid2);
            sleep(10);
            printf("Child1: pid = %d, woke up, exiting\n", getpid());
            return 0;
        }
    } else {
        // ---- original parent ----
        int status;
        waitpid(pid1, &status, 0);
        printf("Original parent: pid = %d, child1 (pid %d) exited\n", getpid(), pid1);
    }

    return 0;
}
\end{lstlisting}

\outputimage{06-recursive-fork.png}

\newpage
\section*{7. Fork and Exec}
\textbf{Question:} Create a child using \texttt{fork()} and replace the
child program using \texttt{execl()}.

\textbf{Code -- fork\_exec.c:}
\begin{lstlisting}
#include <stdio.h>
#include <unistd.h>
#include <sys/wait.h>

int main() {
    printf("fork_exec: parent pid = %d\n", getpid());
    fflush(stdout);

    pid_t pid = fork();

    if (pid < 0) {
        perror("fork failed");
        return 1;
    } else if (pid == 0) {
        // Child process: about to become exec_prog
        printf("fork_exec: child pid = %d, about to exec\n", getpid());
        fflush(stdout);

        execl("./exec_prog", "exec_prog", NULL);

        // only reached if execl fails
        perror("execl failed");
        return 1;
    } else {
        // Parent waits for the child (now running as exec_prog)
        int status;
        waitpid(pid, &status, 0);
        printf("fork_exec: child (pid %d) exited\n", pid);
    }

    return 0;
}
\end{lstlisting}

\textbf{Code -- exec\_program.c:}
\begin{lstlisting}
#include <stdio.h>
#include <unistd.h>

int main() {
    printf("exec_prog: pid = %d, starting count\n", getpid());
    fflush(stdout);

    for (int i = 1; i <= 10; i++) {
        printf("%d\n", i);
        fflush(stdout);
        sleep(1);  // slows it down so you have time to inspect it while running
    }

    return 0;
}
\end{lstlisting}

\outputimage{07-fork-exec.png}

\end{document}
